\chapter{Arrhythmia}
\index{Arrhythmia}

\section*{Definition and Overview}
An arrhythmia is any abnormality of the heart's rhythm: the heartbeat may be too fast (tachycardia), too slow (bradycardia), or irregular. Arrhythmias range from harmless premature beats, which almost everyone experiences at some point, to life-threatening disorders such as ventricular fibrillation. Many arrhythmias produce no symptoms and are detected incidentally, while others cause palpitations, dizziness, or fainting and require treatment. Some, including sustained ventricular tachycardia and severe bradycardia, are medical emergencies. The heart's electrical system normally generates a regular sequence of impulses from the sinoatrial node, and arrhythmias arise when this generation, conduction, or recovery is disturbed.

\section*{Epidemiology}
Arrhythmias are very common and increase with age. Atrial fibrillation, the most common sustained significant arrhythmia, affects around one in ten people over the age of 80 and is a major cause of stroke. Premature atrial and ventricular beats are almost universal and are usually benign. Some inherited arrhythmias, such as long QT syndrome and Brugada syndrome, are rare but can cause sudden cardiac death in young, apparently healthy individuals. Arrhythmias are more common in men and in people with ischaemic heart disease, hypertension, diabetes, heart failure, and obesity. Globally, the burden of arrhythmias is rising in parallel with the ageing population and the prevalence of cardiovascular risk factors.

\section*{Aetiology and Pathophysiology}
An arrhythmia develops when the electrical impulse is generated abnormally, conducts along an abnormal pathway, or is blocked. The mechanism may be an automatic focus firing spontaneously, a re-entry circuit in which the impulse travels in a loop, or impaired conduction between the atria and ventricles. The underlying cause varies.

\begin{itemize}
    \item \textbf{Heart disease:} coronary artery disease, previous myocardial infarction, heart failure, valve disease, or cardiomyopathy
    \item \textbf{Hypertension:} strains and thickens the ventricular muscle over time, favouring arrhythmia
    \item \textbf{Thyroid disease:} hyperthyroidism speeds the heart and predisposes to atrial fibrillation
    \item \textbf{Electrolyte imbalance:} abnormal levels of potassium, calcium, or magnesium disturb the electrical activity of myocytes
    \item \textbf{Stimulants:} caffeine, alcohol, tobacco, and some illicit or recreational drugs
    \item \textbf{Medicines:} certain cold remedies, asthma medications, and some cardiac drugs can provoke arrhythmia
    \item \textbf{Sleep apnoea and stress:} both increase the risk of atrial fibrillation through sympathetic activation and hypoxia
    \item \textbf{Inherited conditions:} channelopathies such as long QT syndrome and Brugada syndrome arise from genetic mutations in cardiac ion channels
    \item \textbf{Ageing:} fibrotic change within the conduction system makes arrhythmia more likely with advancing years
\end{itemize}

\section*{Clinical Features}
Some people with arrhythmias have no symptoms at all, while others feel profoundly unwell. The features depend on the rate, regularity, and duration of the rhythm disturbance and on the function of the underlying heart.

\begin{itemize}
    \item Palpitations: a sensation of fluttering, racing, pounding, or skipped beats
    \item Dizziness, light-headedness, or a feeling of impending faint
    \item Fainting (syncope), sometimes occurring with no warning
    \item Chest discomfort or breathlessness
    \item Fatigue and reduced exercise tolerance
    \item Shortness of breath on exertion or at rest
    \item A sensation that the heart briefly stops or misses a beat
    \item In severe arrhythmias: collapse, loss of pulse, and cardiac arrest
\end{itemize}

\section*{Differential Diagnosis}
The symptoms of arrhythmia can be mimicked by a range of non-cardiac conditions, which should be considered during assessment.

\begin{itemize}
    \item Anxiety, panic attacks, and stress-related symptoms
    \item Anaemia, which can cause palpitations and breathlessness
    \item Hyperthyroidism, with tachycardia, tremor, and heat intolerance
    \item Fever, dehydration, or hypotension from other causes
    \item Overuse of stimulants such as caffeine or energy drinks
    \item Structural problems such as mitral valve prolapse
    \item Non-cardiac chest pain from lung, muscular, or digestive causes
    \item Vasovagal syncope and other causes of transient loss of consciousness, including migraine
\end{itemize}

\section*{When to Seek Medical Care}
See a doctor if palpitations are frequent, prolonged, or accompanied by dizziness, breathlessness, or chest discomfort, or if you have fainted. A family history of arrhythmia or sudden cardiac death should also prompt discussion with a health professional, because some dangerous rhythms are inherited.

\textbf{Contact your local emergency services for:}
\begin{itemize}
    \item Fainting or near-fainting, especially with a rapid or irregular pulse
    \item Chest pain, severe breathlessness, or a racing heart that will not settle
    \item Sudden weakness, confusion, or difficulty speaking
    \item Collapse with unresponsiveness: begin CPR if trained and call emergency services immediately
\end{itemize}

\section*{Diagnosis and Investigations}
The goal of investigation is to capture the heart's rhythm at the moment symptoms occur and to identify any underlying structural or metabolic heart disease.

\begin{itemize}
    \item \textbf{History and examination:} details of symptom onset, triggers, and duration, together with assessment of the pulse, blood pressure, and heart sounds
    \item \textbf{Resting ECG:} a rapid recording of the electrical activity that may show the arrhythmia or its substrate, such as an old infarct
    \item \textbf{Holter monitor:} a portable recorder worn for 24 to 72 hours to capture intermittent rhythms
    \item \textbf{Event or loop recorder:} devices worn or implanted for weeks to months to capture infrequent symptoms
    \item \textbf{Exercise stress test:} records the rhythm during exertion and can expose rate-related abnormalities
    \item \textbf{Echocardiogram:} ultrasound assessment of heart structure, valves, and pump function
    \item \textbf{Blood tests:} electrolytes, thyroid function, renal function, and cardiac enzymes
    \item \textbf{Electrophysiological study:} a specialist catheter test that maps the heart's electrical pathways and can guide ablation
    \item \textbf{Genetic testing:} offered to selected families with suspected inherited arrhythmias
\end{itemize}

\section*{Management}
Treatment depends on the type of arrhythmia, its cause, and the risk it poses. Many benign arrhythmias require no treatment beyond reassurance and avoidance of triggers.

\begin{itemize}
    \item \textbf{Addressing triggers:} reducing caffeine, alcohol, and stress, and treating thyroid disease or sleep apnoea
    \item \textbf{Rate-control medicines:} beta-blockers or calcium-channel blockers to slow a fast ventricular rate
    \item \textbf{Rhythm-control medicines:} anti-arrhythmic drugs such as flecainide or amiodarone to maintain sinus rhythm
    \item \textbf{Anticoagulation:} blood-thinning medicines to reduce stroke risk in atrial fibrillation, guided by risk scores and bleeding risk
    \item \textbf{Cardioversion:} an electric shock or a medicine given to reset the heart to a normal rhythm
    \item \textbf{Catheter ablation:} a procedure that destroys the small area of abnormal tissue generating or perpetuating the arrhythmia
    \item \textbf{Pacemaker:} for symptomatic bradycardia or heart block
    \item \textbf{Implantable defibrillator (ICD):} a device that terminates life-threatening fast rhythms and reduces sudden death risk
    \item \textbf{Follow-up:} regular review with a cardiologist and monitoring of symptoms, devices, and medicines
\end{itemize}

\section*{Complications}
\begin{itemize}
    \item Stroke and other systemic emboli, particularly in atrial fibrillation
    \item Heart failure from a persistently fast or irregular ventricular rate
    \item Tachycardia-induced cardiomyopathy, a weakening of the heart muscle caused by prolonged rapid rhythm
    \item Sudden cardiac arrest in the most dangerous ventricular rhythms
    \item Complications of treatment, including bleeding from anticoagulants and risks associated with ablation or device surgery
    \item Anxiety and reduced quality of life related to worrying symptoms
\end{itemize}

\section*{Prognosis and Outcomes}
The outlook depends strongly on the type of arrhythmia and the presence of underlying heart disease. Many people with benign premature beats or mild arrhythmias have a normal life expectancy. Atrial fibrillation is usually well controlled with rate control, rhythm control, and anticoagulation, but it requires lifelong monitoring. Dangerous ventricular rhythms are serious, yet ablation or an implanted defibrillator can substantially improve survival. Any episode of unexplained fainting or a family history of sudden death warrants investigation, because some causes are treatable and, left unrecognised, can be fatal.

\section*{Prevention and Self-Care}
\begin{itemize}
    \item Keep blood pressure, cholesterol, blood sugar, and weight within a healthy range
    \item Take regular physical activity as advised by your doctor
    \item Limit caffeine and alcohol, and avoid tobacco and stimulant drugs
    \item Manage stress with relaxation, breathing, or mindfulness techniques
    \item Treat contributing conditions such as sleep apnoea and thyroid disease
    \item Take heart medicines exactly as prescribed and never stop them suddenly
    \item If you have a pacemaker or ICD, follow the safety and monitoring instructions given by your team
    \item Know the signs of a cardiac emergency and seek help quickly
\end{itemize}

\section*{Sources and Further Reading}
The following organisations provide reliable, evidence-based information on arrhythmias for patients and healthcare students.

\begin{itemize}
    \item NHS. \textit{Arrhythmia: types, symptoms, and treatment.} https://www.nhs.uk/conditions/arrhythmia/
    \item British Heart Foundation. \textit{Arrhythmias explained.} https://www.bhf.org.uk/informationsupport/heart-matters-hub/medical/arrhythmias
    \item Mayo Clinic. \textit{Heart arrhythmia overview.} https://www.mayoclinic.org/diseases-conditions/heart-arrhythmia/
    \item MedlinePlus. \textit{Arrhythmia.} https://medlineplus.gov/arrhythmia.html
    \item MSD Manual. \textit{Overview of arrhythmias.} https://www.msdmanuals.com/professional/cardiovascular-disorders/arrhythmias-and-conduction-disorders/overview-of-arrhythmias
    \item World Health Organization. \textit{Cardiovascular diseases fact sheet.} https://www.who.int/health-topics/cardiovascular-diseases
\end{itemize}
